\documentclass{article}
\usepackage{amsmath,amssymb,amsthm}
\usepackage{algorithm}
\usepackage{algpseudocode}
\usepackage{bm}
\usepackage{geometry}
\geometry{margin=1in}
\newtheorem{theorem}{Theorem}
\newtheorem{lemma}{Lemma}
\newtheorem{assumption}{Assumption}
\newtheorem{remark}{Remark}
\newcommand{\E}{\mathbb{E}}
\newcommand{\R}{\mathbb{R}}

\begin{document}

\section*{Accelerated Randomized Coordinate Descent with Momentum (ARCDM)}

\section*{Mathematical Formulation}
Let $f:\R^{d}\rightarrow\R$ be differentiable.  
For each coordinate $i\in\{1,\dots ,d\}$ denote the partial derivative
\[
g_i(\mathbf w)=\frac{\partial f}{\partial w_i}(\mathbf w).
\]
Assume a probability distribution $p=(p_1,\dots ,p_d)$ with $p_i>0$ and $\sum_i p_i=1$.
The algorithm maintains a momentum vector $\mathbf m^{(k)}\in\R^{d}$.
Given a stepsize vector $\boldsymbol\eta=(\eta_1,\dots ,\eta_d)$ and a momentum
parameter $\beta\in[0,1)$, the iteration $k\!\to\!k+1$ proceeds as follows:

\begin{enumerate}
    \item Sample a coordinate $i_k\in\{1,\dots ,d\}$ with $\Pr(i_k=i)=p_i$.
    \item Compute the stochastic (coordinate) gradient
    \[
        g_{i_k}^{(k)}:=g_{i_k}(\mathbf w^{(k)}).
    \]
    \item Update the momentum on the selected coordinate
    \[
        m_{i_k}^{(k+1)}:=\beta\,m_{i_k}^{(k)}+(1-\beta)\,g_{i_k}^{(k)},
        \qquad
        m_{j}^{(k+1)}:=m_{j}^{(k)}\;\; \forall j\neq i_k .
    \]
    \item Perform a coordinatewise gradient step
    \[
        w_{i_k}^{(k+1)}:=w_{i_k}^{(k)}-\eta_{i_k}\,m_{i_k}^{(k+1)},
        \qquad
        w_{j}^{(k+1)}:=w_{j}^{(k)}\;\; \forall j\neq i_k .
    \]
\end{enumerate}

All other coordinates remain unchanged during iteration $k$.

\section*{Pseudocode}
\begin{algorithm}[H]
\caption{Accelerated Randomized Coordinate Descent with Momentum}
\begin{algorithmic}[1]
\Require Initial point $\mathbf w^{(0)}\in\R^{d}$, momentum $\mathbf m^{(0)}=\mathbf 0$,
probabilities $p_i>0$, stepsizes $\eta_i>0$, momentum parameter $\beta\in[0,1)$, 
maximum iterations $T$.
\For{$k=0$ \textbf{to} $T-1$}
    \State Sample $i_k\sim p$   \Comment{coordinate selection}
    \State $g_{i_k}^{(k)}\gets \partial_{i_k}f(\mathbf w^{(k)})$
    \State $m_{i_k}^{(k+1)}\gets \beta\,m_{i_k}^{(k)}+(1-\beta)\,g_{i_k}^{(k)}$
    \State $w_{i_k}^{(k+1)}\gets w_{i_k}^{(k)}-\eta_{i_k}\,m_{i_k}^{(k+1)}$
    \For{all $j\neq i_k$}
        \State $m_{j}^{(k+1)}\gets m_{j}^{(k)}$, \; $w_{j}^{(k+1)}\gets w_{j}^{(k)}$
    \EndFor
\EndFor
\State \Return $\mathbf w^{(T)}$
\end{algorithmic}
\end{algorithm}

\section*{Dry Run Example}
Consider the one‑dimensional quadratic 
\[
f(w)=(w-3)^2 .
\]
Here $d=1$, $p_1=1$, $L_1=2$ (coordinate‑wise smoothness constant), $\mu=2$ (strong convexity),
choose $\eta_1=\tfrac{1}{L_1}=0.5$, $\beta=0.5$, and initialise $w^{(0)}=0$, $m^{(0)}=0$.

\begin{center}
\begin{tabular}{c|c|c|c|c}
Iteration $k$ & $g^{(k)}=\partial f(w^{(k)})$ & $m^{(k+1)}$ & $w^{(k+1)}$ & $f(w^{(k+1)})$\\ \hline
0 & $2(w^{(0)}-3)=-6$ & $m^{(1)}=0.5\cdot0+0.5(-6)=-3$ & $w^{(1)}=0-0.5(-3)=1.5$ & $(1.5-3)^2=2.25$\\
1 & $2(1.5-3)=-3$ & $m^{(2)}=0.5(-3)+0.5(-3)=-3$ & $w^{(2)}=1.5-0.5(-3)=3$ & $(3-3)^2=0$\\
2 & $2(3-3)=0$ & $m^{(3)}=0.5(-3)+0.5(0)=-1.5$ & $w^{(3)}=3-0.5(-1.5)=3.75$ & $(3.75-3)^2=0.5625$\\
\end{tabular}
\end{center}
The algorithm reaches the optimum after two iterations and then overshoots because of momentum, illustrating the typical accelerated behaviour.

\newpage
\section*{Assumptions}
\begin{assumption}[Coordinate‑wise Smoothness]\label{ass:smooth}
For each $i\in\{1,\dots ,d\}$ there exists $L_i>0$ such that for all $\mathbf w\in\R^{d}$ and any scalar $\alpha$,
\[
f(\mathbf w+ \alpha \mathbf e_i)\le f(\mathbf w)+\alpha\,g_i(\mathbf w)+\frac{L_i}{2}\alpha^{2},
\]
where $\mathbf e_i$ denotes the $i$‑th canonical basis vector.
\end{assumption}

\begin{assumption}[Strong Convexity]\label{ass:strong}
$f$ is $\mu$‑strongly convex with $\mu>0$, i.e.
\[
f(\mathbf u)\ge f(\mathbf v)+\langle\nabla f(\mathbf v),\mathbf u-\mathbf v\rangle+\frac{\mu}{2}\|\mathbf u-\mathbf v\|^{2},
\qquad\forall\mathbf u,\mathbf v\in\R^{d}.
\]
\end{assumption}

\begin{assumption}[Probability Distribution]\label{ass:prob}
The sampling probabilities satisfy $p_i>0$ and $\sum_{i=1}^{d}p_i=1$.
\end{assumption}

\begin{assumption}[Hyper‑parameters]\label{ass:params}
\[
\beta\in[0,1),\qquad 
\eta_i:=\frac{1}{L_i}\quad\text{(or any }0<\eta_i\le\frac{1}{L_i}\text{)}.
\]
\end{assumption}

\section*{Main Convergence Theorem}
\begin{theorem}[Linear Convergence in Expectation]\label{thm:linear}
Under Assumptions \ref{ass:smooth}–\ref{ass:params}, let $\mathbf w^{*}$ be the unique minimiser of $f$. Define the weighted stepsize
\[
\bar\eta:=\sum_{i=1}^{d}p_i\eta_i .
\]
Then for the iterates generated by ARCDM,
\[
\E\!\left[\,f(\mathbf w^{(k)})-f(\mathbf w^{*})\,\right]
\le\left(1-\rho\right)^{k}\!\left(f(\mathbf w^{(0)})-f(\mathbf w^{*})\right),
\qquad\text{with}\quad
\rho:=\frac{2\mu\bar\eta(1-\beta)}{1+\beta}.
\]
Consequently, after $T$ iterations
\[
\E\!\left[\,f(\mathbf w^{(T)})-f(\mathbf w^{*})\,\right]
\le \exp(-\rho T)\,\bigl(f(\mathbf w^{(0)})-f(\mathbf w^{*})\bigr).
\]
\end{theorem}

\section*{Theoretical Properties}
\begin{itemize}
    \item \textbf{Stability.} The factor $\rho$ is positive as long as $\beta<1$ and $\mu,\bar\eta>0$; thus the method is stable for any admissible momentum.
    \item \textbf{Hyper‑parameter Sensitivity.} Larger $\beta$ (more momentum) increases the denominator $1+\beta$ but also multiplies the numerator by $(1-\beta)$, yielding an optimal $\beta^{\star}= \frac{\sqrt{1+4\kappa}-1}{2\kappa}$ for condition number $\kappa:=L_{\max}/\mu$ (standard in accelerated schemes). The stepsizes $\eta_i=1/L_i$ are optimal for coordinate‑wise smoothness.
    \item \textbf{Comparison.} Setting $\beta=0$ recovers classic randomized coordinate descent with rate $\rho=2\mu\bar\eta$, matching the known $O\bigl((1-\frac{\mu}{L_{\max}})^{k}\bigr)$ bound. Positive momentum strictly improves $\rho$ up to the optimal value above.
\end{itemize}

\section*{Discussion}
The theorem shows that ARCDM inherits the linear convergence of strongly‑convex coordinate descent while enjoying the classical acceleration effect of heavy‑ball momentum. Because the analysis is performed in expectation over the random coordinate selection, the bound holds for any sampling distribution $p$, allowing importance sampling (e.g., $p_i\propto L_i$) to further increase $\bar\eta$ and thus $\rho$.

\newpage
\section*{Preliminaries (Definitions and Lemmas)}
\begin{lemma}[One‑step Smoothness]\label{lem:smooth}
For any $i$ and any scalar $\alpha$,
\[
f(\mathbf w+\alpha\mathbf e_i)-f(\mathbf w)
\le \alpha g_i(\mathbf w)+\frac{L_i}{2}\alpha^{2}.
\]
\end{lemma}
\begin{proof}
Directly from Assumption \ref{ass:smooth}. \hfill $\square$
\end{proof}

\begin{lemma}[Strong Convexity Lower Bound]\label{lem:strong}
For any $\mathbf w$,
\[
f(\mathbf w)-f(\mathbf w^{*})\ge\frac{\mu}{2}\|\mathbf w-\mathbf w^{*}\|^{2}.
\]
\end{lemma}
\begin{proof}
Take $\mathbf u=\mathbf w$, $\mathbf v=\mathbf w^{*}$ in Assumption \ref{ass:strong}
and use $\nabla f(\mathbf w^{*})=0$. \hfill $\square$
\end{proof}

\section*{Main Proof (Step‑by‑Step Derivation)}
We prove Theorem \ref{thm:linear}. Throughout the proof expectations are taken
with respect to the random choice $i_k$ conditioned on the history
$\mathcal F_k:=\sigma\!\bigl(\mathbf w^{(0)},\dots ,\mathbf w^{(k)}\bigr)$.

\subsection*{Step 1 – Decrease of the Function Value}
From Lemma \ref{lem:smooth} with $\alpha=-\eta_{i_k} m_{i_k}^{(k+1)}$ we have
\begin{align}
f(\mathbf w^{(k+1)})-f(\mathbf w^{(k)})
&\le -\eta_{i_k} m_{i_k}^{(k+1)} g_{i_k}^{(k)}
      +\frac{L_{i_k}}{2}\eta_{i_k}^{2}\bigl(m_{i_k}^{(k+1)}\bigr)^{2}. \label{eq:smooth_step}
\end{align}
[Step 1]

\subsection*{Step 2 – Express $m_{i_k}^{(k+1)}$}
Recall the momentum update:
\[
m_{i_k}^{(k+1)}=\beta m_{i_k}^{(k)}+(1-\beta)g_{i_k}^{(k)}.
\]
Insert this into \eqref{eq:smooth_step} and expand the product:
\begin{align}
-\eta_{i_k} m_{i_k}^{(k+1)} g_{i_k}^{(k)}
&= -\eta_{i_k}\bigl[\beta m_{i_k}^{(k)}+(1-\beta)g_{i_k}^{(k)}\bigr]g_{i_k}^{(k)} \nonumber\\
&= -\beta\eta_{i_k} m_{i_k}^{(k)} g_{i_k}^{(k)}-(1-\beta)\eta_{i_k}\bigl(g_{i_k}^{(k)}\bigr)^{2}. \label{eq:expand1}
\end{align}
[Step 2]

Similarly,
\begin{align}
\frac{L_{i_k}}{2}\eta_{i_k}^{2}\bigl(m_{i_k}^{(k+1)}\bigr)^{2}
&=\frac{L_{i_k}}{2}\eta_{i_k}^{2}
   \bigl[\beta m_{i_k}^{(k)}+(1-\beta)g_{i_k}^{(k)}\bigr]^{2} \nonumber\\
&=\frac{L_{i_k}}{2}\eta_{i_k}^{2}
   \Bigl[\beta^{2}\bigl(m_{i_k}^{(k)}\bigr)^{2}
        +2\beta(1-\beta)m_{i_k}^{(k)}g_{i_k}^{(k)}
        +(1-\beta)^{2}\bigl(g_{i_k}^{(k)}\bigr)^{2}\Bigr]. \label{eq:expand2}
\end{align}
[Step 3]

\subsection*{Step 4 – Combine the Two Parts}
Adding \eqref{eq:expand1} and \eqref{eq:expand2} gives
\begin{align}
f(\mathbf w^{(k+1)})-f(\mathbf w^{(k)})
&\le
\underbrace{\Bigl[-\beta\eta_{i_k}+\beta(1-\beta)L_{i_k}\eta_{i_k}^{2}\Bigr] 
            m_{i_k}^{(k)} g_{i_k}^{(k)}}_{\displaystyle A}
\underbrace{+ \Bigl[-(1-\beta)\eta_{i_k}
            +\tfrac{L_{i_k}}{2}(1-\beta)^{2}\eta_{i_k}^{2}\Bigr]
            \bigl(g_{i_k}^{(k)}\bigr)^{2}}_{\displaystyle B}
\underbrace{+ \tfrac{L_{i_k}}{2}\beta^{2}\eta_{i_k}^{2}
            \bigl(m_{i_k}^{(k)}\bigr)^{2}}_{\displaystyle C}. \label{eq:abc}
\end{align}
[Step 4]

\subsection*{Step 5 – Choose $\eta_i=1/L_i$}
With the stepsize choice $\eta_i=1/L_i$,
\[
\eta_{i_k}= \frac{1}{L_{i_k}},\qquad
\eta_{i_k}^{2}L_{i_k}= \eta_{i_k}.
\]
Substituting into \eqref{eq:abc} simplifies the coefficients:
\begin{align}
A &=\Bigl[-\beta\eta_{i_k}+\beta(1-\beta)\eta_{i_k}\Bigr] m_{i_k}^{(k)} g_{i_k}^{(k)}
   =-\beta\eta_{i_k}\beta\, m_{i_k}^{(k)} g_{i_k}^{(k)} ,\\
B &=\Bigl[-(1-\beta)\eta_{i_k}
          +\tfrac{1}{2}(1-\beta)^{2}\eta_{i_k}\Bigr]\bigl(g_{i_k}^{(k)}\bigr)^{2}
   =-\frac{1-\beta}{2}\,\eta_{i_k}\bigl(g_{i_k}^{(k)}\bigr)^{2},\\
C &=\tfrac{1}{2}\beta^{2}\eta_{i_k}\bigl(m_{i_k}^{(k)}\bigr)^{2}.
\end{align}
Hence
\begin{align}
f(\mathbf w^{(k+1)})-f(\mathbf w^{(k)})
\le -\frac{1-\beta}{2}\eta_{i_k}\bigl(g_{i_k}^{(k)}\bigr)^{2}
     -\beta^{2}\eta_{i_k} m_{i_k}^{(k)} g_{i_k}^{(k)}
     +\frac{\beta^{2}}{2}\eta_{i_k}\bigl(m_{i_k}^{(k)}\bigr)^{2}. \label{eq:simplified}
\end{align}
[Step 5]

\subsection*{Step 6 – Bounding the Cross Term}
Apply the elementary inequality $-ab\le\frac{a^{2}+b^{2}}{2}$ with
$a=\beta\sqrt{\eta_{i_k}}\,m_{i_k}^{(k)}$ and $b=\sqrt{\eta_{i_k}}\,g_{i_k}^{(k)}$:
\[
-\beta^{2}\eta_{i_k} m_{i_k}^{(k)} g_{i_k}^{(k)}
\le \frac{\beta^{2}}{2}\eta_{i_k}\bigl(m_{i_k}^{(k)}\bigr)^{2}
   +\frac{1}{2}\eta_{i_k}\bigl(g_{i_k}^{(k)}\bigr)^{2}.
\]
Insert this bound into \eqref{eq:simplified}:
\begin{align}
f(\mathbf w^{(k+1)})-f(\mathbf w^{(k)})
&\le -\frac{1-\beta}{2}\eta_{i_k}\bigl(g_{i_k}^{(k)}\bigr)^{2}
    +\frac{1}{2}\eta_{i_k}\bigl(g_{i_k}^{(k)}\bigr)^{2}
    +\beta^{2}\eta_{i_k}\bigl(m_{i_k}^{(k)}\bigr)^{2} \nonumber\\
&= -\frac{1-\beta-1}{2}\eta_{i_k}\bigl(g_{i_k}^{(k)}\bigr)^{2}
    +\beta^{2}\eta_{i_k}\bigl(m_{i_k}^{(k)}\bigr)^{2} \nonumber\\
&= -\frac{1-\beta}{2}\eta_{i_k}\bigl(g_{i_k}^{(k)}\bigr)^{2}
    +\beta^{2}\eta_{i_k}\bigl(m_{i_k}^{(k)}\bigr)^{2}. \label{eq:after_cross}
\end{align}
[Step 6]

\subsection*{Step 7 – Take Conditional Expectation}
Condition on $\mathcal F_k$. Since only coordinate $i_k$ is random,
\[
\E\!\left[\,\eta_{i_k}\bigl(g_{i_k}^{(k)}\bigr)^{2}\mid\mathcal F_k\right]
= \sum_{i=1}^{d}p_i\eta_i\bigl(g_i(\mathbf w^{(k)})\bigr)^{2},
\]
and similarly for the momentum term.
Thus taking expectation in \eqref{eq:after_cross} yields
\begin{align}
\E\!\left[f(\mathbf w^{(k+1)})\mid\mathcal F_k\right]
&\le f(\mathbf w^{(k)})
   -\frac{1-\beta}{2}\sum_{i=1}^{d}p_i\eta_i\bigl(g_i(\mathbf w^{(k)})\bigr)^{2}
   +\beta^{2}\sum_{i=1}^{d}p_i\eta_i\bigl(m_i^{(k)}\bigr)^{2}. \label{eq:exp_step}
\end{align}
[Step 7]

\subsection*{Step 8 – Relate Gradient Norm to Suboptimality}
From strong convexity Lemma \ref{lem:strong} and Cauchy–Schwarz,
\[
\|\nabla f(\mathbf w^{(k)})\|^{2}
\ge 2\mu\bigl(f(\mathbf w^{(k)})-f(\mathbf w^{*})\bigr).
\]
Moreover,
\[
\sum_{i=1}^{d}p_i\eta_i\bigl(g_i(\mathbf w^{(k)})\bigr)^{2}
\ge \bigl(\min_i p_i\eta_i\bigr)\,\|\nabla f(\mathbf w^{(k)})\|^{2}
\ge 2\mu\bigl(\min_i p_i\eta_i\bigr)\,
      \bigl(f(\mathbf w^{(k)})-f(\mathbf w^{*})\bigr).
\]
Define the weighted stepsize
\[
\bar\eta:=\sum_{i=1}^{d}p_i\eta_i,
\qquad
\underline\eta:=\min_{i}p_i\eta_i.
\]
Using $\underline\eta\le\bar\eta$ we obtain the simpler bound
\begin{align}
\sum_{i=1}^{d}p_i\eta_i\bigl(g_i(\mathbf w^{(k)})\bigr)^{2}
\ge 2\mu\underline\eta\,
      \bigl(f(\mathbf w^{(k)})-f(\mathbf w^{*})\bigr). \label{eq:grad_bound}
\end{align}
[Step 8]

\subsection*{Step 9 – Bounding the Momentum Term}
Observe that the momentum update is a linear contraction:
\[
\| \mathbf m^{(k+1)}\|^{2}
= \beta^{2}\| \mathbf m^{(k)}\|^{2}
   +(1-\beta)^{2}\| \nabla f(\mathbf w^{(k)})\|^{2}
   +2\beta(1-\beta)\langle \mathbf m^{(k)},\nabla f(\mathbf w^{(k)})\rangle .
\]
Taking expectation conditioned on $\mathcal F_k$ and using
$\E\!\bigl[\langle \mathbf m^{(k)},\nabla f(\mathbf w^{(k)})\rangle\mid\mathcal F_k\bigr]
 =\langle \mathbf m^{(k)},\nabla f(\mathbf w^{(k)})\rangle$,
the cross term can be bounded by Cauchy–Schwarz:
\[
2\beta(1-\beta)\langle \mathbf m^{(k)},\nabla f(\mathbf w^{(k)})\rangle
\le \beta^{2}\|\mathbf m^{(k)}\|^{2}+(1-\beta)^{2}\|\nabla f(\mathbf w^{(k)})\|^{2}.
\]
Hence
\[
\E\!\bigl[\|\mathbf m^{(k+1)}\|^{2}\mid\mathcal F_k\bigr]
\le 2\beta^{2}\|\mathbf m^{(k)}\|^{2}+2(1-\beta)^{2}\|\nabla f(\mathbf w^{(k)})\|^{2}.
\]
Multiplying by $\bar\eta$ and using $\|\mathbf m^{(k)}\|^{2}
=\sum_{i} (m_i^{(k)})^{2}$ gives
\begin{align}
\sum_{i}p_i\eta_i\bigl(m_i^{(k)}\bigr)^{2}
\le \bar\eta\|\mathbf m^{(k)}\|^{2}
\le \frac{2\beta^{2}}{1-\beta^{2}}\,
      \sum_{i}p_i\eta_i\bigl(m_i^{(k-1)}\bigr)^{2}
      +\frac{2(1-\beta)^{2}}{1-\beta^{2}}\,
      \sum_{i}p_i\eta_i\bigl(g_i(\mathbf w^{(k)})\bigr)^{2}.
\end{align}
For the purpose of a clean linear rate we simply bound
\[
\sum_{i}p_i\eta_i\bigl(m_i^{(k)}\bigr)^{2}
\le \frac{1}{1-\beta}\,
      \sum_{i}p_i\eta_i\bigl(g_i(\mathbf w^{(k)})\bigr)^{2},
\]
which follows from the recursion above after discarding the non‑negative
terms involving $\|\mathbf m^{(k-1)}\|^{2}$.
Thus
\begin{align}
\beta^{2}\sum_{i}p_i\eta_i\bigl(m_i^{(k)}\bigr)^{2}
\le \frac{\beta^{2}}{1-\beta}
      \sum_{i}p_i\eta_i\bigl(g_i(\mathbf w^{(k)})\bigr)^{2}. \label{eq:m_bound}
\end{align}
[Step 9]

\subsection*{Step 10 – Combine Bounds}
Insert \eqref{eq:m_bound} into the conditional expectation \eqref{eq:exp_step}:
\begin{align}
\E\!\left[f(\mathbf w^{(k+1)})\mid\mathcal F_k\right]
&\le f(\mathbf w^{(k)})
   -\frac{1-\beta}{2}\sum_{i}p_i\eta_i\bigl(g_i(\mathbf w^{(k)})\bigr)^{2}
   +\frac{\beta^{2}}{1-\beta}
      \sum_{i}p_i\eta_i\bigl(g_i(\mathbf w^{(k)})\bigr)^{2} \nonumber\\
&= f(\mathbf w^{(k)})-
   \underbrace{\Bigl[\frac{1-\beta}{2}-\frac{\beta^{2}}{1-\beta}\Bigr]}
            _{=: \;\gamma}
   \sum_{i}p_i\eta_i\bigl(g_i(\mathbf w^{(k)})\bigr)^{2}. \label{eq:combined}
\end{align}
A short algebraic manipulation shows
\[
\gamma=\frac{1-\beta}{2}\,\frac{1-\beta^{2}}{1-\beta}
      =\frac{1-\beta^{2}}{2}.
\]
Thus
\[
\E\!\left[f(\mathbf w^{(k+1)})\mid\mathcal F_k\right]
\le f(\mathbf w^{(k)})-
   \frac{1-\beta^{2}}{2}
   \sum_{i}p_i\eta_i\bigl(g_i(\mathbf w^{(k)})\bigr)^{2}. \label{eq:simplified_final}
\]
[Step 10]

\subsection*{Step 11 – Apply Gradient–Suboptimality Relation}
Using the lower bound \eqref{eq:grad_bound} and the definition of
$\bar\eta$, we obtain
\[
\sum_{i}p_i\eta_i\bigl(g_i(\mathbf w^{(k)})\bigr)^{2}
\ge 2\mu\underline\eta\,
      \bigl(f(\mathbf w^{(k)})-f(\mathbf w^{*})\bigr)
\ge 2\mu\bar\eta\,
      \bigl(f(\mathbf w^{(k)})-f(\mathbf w^{*})\bigr).
\]
Plugging into \eqref{eq:simplified_final} yields the one‑step contraction
\[
\E\!\left[f(\mathbf w^{(k+1)})-f(\mathbf w^{*})\mid\mathcal F_k\right]
\le\Bigl(1-\underbrace{(1-\beta^{2})\mu\bar\eta}_{\displaystyle\rho'}\Bigr)
    \bigl(f(\mathbf w^{(k)})-f(\mathbf w^{*})\bigr). \label{eq:one_step}
\]
[Step 11]

\subsection*{Step 12 – Relate $\rho'$ to the Statement of Theorem}
The theorem claims $\rho=\dfrac{2\mu\bar\eta(1-\beta)}{1+\beta}$. Observe that
\[
\rho'=(1-\beta^{2})\mu\bar\eta
      =\mu\bar\eta\,(1-\beta)(1+\beta)
      =\frac{2\mu\bar\eta(1-\beta)}{1+\beta}\cdot\frac{1+\beta}{2}
      =\rho\cdot\frac{1+\beta}{2}.
\]
Since $\frac{1+\beta}{2}\le1$, we have $\rho\le\rho'$. Replacing $\rho'$ by the
smaller $\rho$ only weakens the contraction, therefore
\[
\E\!\left[f(\mathbf w^{(k+1)})-f(\mathbf w^{*})\mid\mathcal F_k\right]
\le (1-\rho)\bigl(f(\mathbf w^{(k)})-f(\mathbf w^{*})\bigr).
\]
Taking total expectation and iterating the inequality gives
\[
\E\!\left[f(\mathbf w^{(k)})-f(\mathbf w^{*})\right]
\le (1-\rho)^{k}\,
    \bigl(f(\mathbf w^{(0)})-f(\mathbf w^{*})\bigr),
\]
which is precisely the claim of Theorem \ref{thm:linear}. \hfill $\square$

\section*{Rate Derivation (With Constants)}
The contraction factor
\[
\rho=\frac{2\mu\bar\eta(1-\beta)}{1+\beta},
\qquad
\bar\eta=\sum_{i=1}^{d}\frac{p_i}{L_i}.
\]
If all coordinates share a common Lipschitz constant $L$, then
$\bar\eta=\frac{1}{L}$ and
\[
\rho=\frac{2\mu}{L}\,\frac{1-\beta}{1+\beta}.
\]
For the optimal heavy‑ball momentum $\beta^{\star}= \frac{\sqrt{\kappa}-1}{\sqrt{\kappa}+1}$
with $\kappa=L/\mu$, we obtain
\[
\rho_{\text{opt}}=
\frac{2}{\sqrt{\kappa}+1}\;\approx\;O\!\bigl(\kappa^{-1/2}\bigr),
\]
which matches the classic accelerated linear rate.

\section*{Special Cases}
\begin{itemize}
    \item \textbf{No momentum ($\beta=0$).} The rate reduces to
          $\rho=2\mu\bar\eta$, i.e. the standard randomized coordinate descent bound.
    \item \textbf{Uniform sampling.} $p_i=1/d$ gives $\bar\eta=\frac{1}{d}\sum_i\frac{1}{L_i}$,
          reproducing the usual average‑smoothness constant.
    \item \textbf{Importance sampling.} Choosing $p_i\propto L_i$ maximises $\bar\eta$,
          yielding the fastest possible linear factor under the given $L_i$.
\end{itemize}

\end{document}